\documentclass[12pt]{article}
\usepackage{geometry}
\usepackage{booktabs}
\usepackage{longtable}
\usepackage{array}
\usepackage{hyperref}
\usepackage{enumitem}
\geometry{margin=1in}
\title{Completed Research Portfolio Summary}
\author{JC Vaught}
\date{October 2, 2025}
\begin{document}
\maketitle
\section*{Scope}
This report consolidates completed research deliverables stored within the project workspace as of October~2,~2025. Each section captures outcomes that have already been produced, verified, or documented; prospective plans and future-facing proposals that appear in the source materials have been intentionally excluded in accordance with the reporting requirement.

\section{GPU Parallelism and Scaling Studies}
\subsection{Literature Synthesis}
A comparative literature survey distilled the canonical sharding strategies for transformer-scale training across data, model, pipeline, and tensor parallelism. The deliverable summarizes implementation complexity, strengths, and latency/communication trade-offs to guide deployment choices on small GPU clusters.\cite{gpu-summary-table}

\subsection{Benchmarked Workflows}
Hands-on experiments were executed across two- and four-GPU configurations using transformer language models at 1{,}024- and 4{,}096-dimension embeddings. The study logged throughput, memory ceilings, utilization, and recovery characteristics for each parallelism mode, providing concrete guidance on when the added orchestration cost delivers a measurable benefit.\cite{gpu-weekly-brief}

\section{Thermal Camera Calibration Survey}
A comprehensive review of non-radar calibration targets categorized active versus passive thermal patterns and captured achieved reprojection errors where available.\cite{thermal-summary} Table~\ref{tab:thermal-calibration} enumerates the documented configurations.

\begin{longtable}{>{\raggedright\arraybackslash}p{3cm} >{\raggedright\arraybackslash}p{2cm} >{\raggedright\arraybackslash}p{2.5cm} >{\raggedright\arraybackslash}p{6cm}}
\caption{Completed thermal calibration approaches and performance metrics.}\label{tab:thermal-calibration}\\
\toprule
Method & Category & RMSE (px) & Notes \\
\midrule
Paper checkerboard + flood lamp & Passive & 0.41--0.48 & Heated board delivers sub-half-pixel accuracy with minimal equipment. \\
IR lamp checkerboard (HDR frame blend) & Passive & 0.60--0.714 & Increased dynamic range required to visualize pattern, but error increases because of lamp non-uniformity. \\
Monitor-displayed mask & Passive & 0.284--0.324 & Utilizes emissive screen for high-contrast target without custom fabrication. \\
Aluminum tape on board & Passive & --- & Commonly deployed low-cost method; qualitative assessment notes strong contrast from metal's low emissivity. \\
Printed aluminum plate (ink) & Passive & 0.10--0.22 & Metal substrate with printed pattern delivers best-in-class passive accuracy; commercially common. \\
Printed aluminum plate (UV) & Passive & 0.233 & UV-cured print variant yields marginally higher error than inked plate. \\
Heated resistive wires & Active & 0.175--0.177 & Precisely positioned heat points create reliable features for calibration. \\
Bulb array on backing & Active & 0.588 & Bulb heating provides workable calibration but suffers from thermal bloom. \\
Circular mask target & Passive & 0.052 & Circular aperture mask documented as the lowest error among surveyed methods. \\
\bottomrule
\end{longtable}

Cost analysis highlighted metal targets (e.g., 24\,×\,24 in plates quoted at \$55) and active heating rigs as the primary expense drivers, motivating preference for reusable passive configurations with demonstrated sub-0.2 pixel accuracy.\cite{thermal-summary}

\section{Maritime Vision Dataset Catalog}
The dataset catalog consolidates public maritime sensing corpora across electro-optical, infrared, and radar modalities to accelerate training data acquisition for fusion pipelines.\cite{maritime-datasets} Key exemplars include:
\begin{itemize}[leftmargin=*,nosep]
  \item Singapore Maritime Dataset (SMD): 81 on-shore/on-board 1080p video sequences with day/night/haze diversity for detection and tracking benchmarks.
  \item MARVEL: \textasciitilde2~million web-crawled ship images across 109 vessel classes for fine-grained classification.
  \item MarDCT (ARGOS-Venice): 6{,}742 annotated images plus 16 Venice waterway videos supporting urban USV research.
  \item SeaDronesSee: >54{,}000 UAV frames with \textasciitilde400k annotations plus multispectral bands for search-and-rescue detection.
  \item KOLOMVERSE: 186{,}419 4K images extracted from 5{,}845 hours of CCTV coverage on 17 vessels across 21 sea zones, with bounding boxes for collision avoidance.
  \item M2SODAI: Multi-modal aerial RGB + hyperspectral dataset capturing 5{,}764 ship and debris instances to suppress sun-glint false alarms.
\end{itemize}
The catalog couples each entry with publication references, sensor modalities, and download provenance, easing compliance and reproducibility tracking for downstream experiments.

\section{Radar--Camera Fusion Research}
A technical whitepaper documented the empirical limitations of X-band radar and optical-only detection for small, low-RCS targets in dynamic sea clutter, then proposed a radar-assisted annotation pipeline.\cite{fusion-summary} Completed analyses include:
\begin{itemize}[leftmargin=*,nosep]
  \item Demonstrated CFAR thresholding on Lake Murray radar captures to illustrate clutter-adaptive detection and its residual failure modes in inhomogeneous seas.
  \item Quantified automatic gain control side effects that either suppress or saturate small-object returns, establishing the requirement for raw-data access and multi-gain capture.
  \item Outlined a co-registration workflow where radar-projected regions-of-interest accelerate human-in-the-loop labeling, reducing manual scan burden for sub-5~pixel targets.
\end{itemize}

\section{M364C FLIR Sensor Calibration and Data Pipeline (May 2025)}
The ONR monthly report for May~2025 compiled finished deliverables across imaging, fusion, and range estimation tasks.\cite{monthly-report-may2025}
\subsection{Imaging System Enhancements}
\begin{itemize}[leftmargin=*,nosep]
  \item Upgraded the Jetson Orin camera GUI with real-time uptime, CPU/RAM/temperature telemetry, storage endurance estimates, selectable overlays, and full-screen per-camera views (Figure~1 in source).
  \item Fabricated a 60~mm aluminum-square thermal calibration board, enabling reliable IR visualization of checkerboard corners for the M364C.
  \item Processed paired IR/RGB captures through MATLAB Stereo Calibrator, achieving 0.80~pixel mean reprojection error.
  \item Explored uncalibrated SURF-based rectification, using pixel-difference diagnostics to evaluate residual misalignment between modalities.
  \item Captured new Charleston, SC shoreline sequences, expanding the IMSEL maritime dataset to 8{,}589 annotated IR/RGB pairs (1{,}411 short of the 10k goal).
\end{itemize}

\subsection{Geometric Processing and Data Management}
\begin{itemize}[leftmargin=*,nosep]
  \item Replaced fixed crop windows with a world-distance cutoff to define rectified image bounds, keeping edge-length errors under 10\% on tiled ground truth scenes.
  \item Commissioned a radar/camera/INS recorder on Jetson AGX Orin; timestamped sweeps, frames, and navigation metadata now persist without loss for downstream fusion alignment.
\end{itemize}

\subsection{Range Estimation and Environmental Modeling}
\begin{itemize}[leftmargin=*,nosep]
  \item Refined geometric range calibration by acquiring camera--IMU transform matrices at \textasciitilde10~m board separation and logging sensor time shifts.
  \item Improved LWIR range regression by reparameterizing inputs to XGBoost, lifting $R^2$ from 0.65 to 0.90 across seven datasets spanning 14 unknown gain settings; captured 250+ paired LWIR frames with ground-truth distances using a demo Axiom Optics sensor.
  \item Trained a rain segmentation model and instituted manual-cleanup workflows to raise mask fidelity for deraining efforts.
\end{itemize}

\section{Multi-Gain Radar Characterization (June 2025)}
A focused campaign partially validated the multi-gain acquisition workflow and streamlined the Furuno NXT capture stack.\cite{phase2-gain}
\begin{itemize}[leftmargin=*,nosep]
  \item Gathered \textasciitilde60{,}000 sweeps (\textasciitilde200~GB) across seven Lake Murray and Lake Greenwood access points over 20 days, covering day, night, and storm conditions with 1.5~NM effective radius per site.
  \item Entropy analysis over five ranges and 100 gain steps (5 samples each) identified the 30--93\% S1/S2 zone as information-rich while flagging saturation above 94\%.
  \item Logged 2{,}500 sweeps per mini-calibration in 40~minutes, informing design of a 20k-sweep comprehensive calibration.
  \item Hardened field hardware with a new enclosure and OS tuning, shrinking Raspberry Pi boot time from $\ge$4~min to $<20$~s.
  \item Authored a parallel CSV-to-NetCDF converter that compresses storage by \textasciitilde25\times and attaches richer metadata; multi-core implementation accelerates cleanup by $>$4\times relative to previous scripts.
  \item Benchmarked six single-frame denoisers (BM3D best) and multi-frame filters (multi-look median leading) plus OpenCV HDR merges (Merge Mertens and NLM improving low-SNR targets).
  \item Trained an initial temporal masked autoencoder on 132 samples (40 epochs, 75\% masking) to assess feasibility of self-supervised radar reconstruction.
\end{itemize}

\section{HDR Radar Tracking and Simulation (July 2025)}
Subsequent work operationalized HDR fusion and automated tracking on newly collected datasets.\cite{phase2-hdr}
\begin{itemize}[leftmargin=*,nosep]
  \item Characterized receiver response via dual corner reflectors, noting compression above gain level 9 (0--100 scale) and capturing parameter scripts for repeatability.
  \item Selected gain triplet (40/50/70) and short-2 pulse for lake operations; Robertson HDR fusion produced 16-bit PPIs, boosting peak signal-to-clutter by $>$12~dB and visibly clarifying land, buoys, and boats.
  \item Built a Monte Carlo simulator that ingests GIS shorelines and buoy catalogs, supports Rayleigh/K/Weibull clutter models, allows STC curve sweeps, and matches real scenes within 2~dB across range gates.
  \item Delivered two functional trackers: a self-initializing extended Kalman filter with adaptive birth/dormant logic, and a temporal masked autoencoder fine-tuned on three hand-labelled masks that already exceeds July's baseline IoU.
\end{itemize}

\section{Site Surveys for Radar Data Collection}
The South Carolina lake access workbook catalogs actionable launch points, geocodes, and amenity notes for inland and coastal radar deployments.\cite{lake-access}
\begin{itemize}[leftmargin=*,nosep]
  \item Compiled access documentation for Lake Marion, Moultrie, Murray, Hartwell, Keowee, Jocassee, and Thurmond, each with links to operating authorities and ramp directories.
  \item Listed Google Maps references and coordinates for Monticello access points (e.g., 34.324162~N, 81.285518~W) to streamline logistics for night operations.
  \item Extended coverage to coastal ramps (Calibogue, Port Royal, St.~Helena Sounds) and priority beach sites, capturing facility notes such as restroom and shower availability relevant to crew planning.
\end{itemize}

\section{Portable Power Station Comparison}
A market survey benchmarked eight sub-350~Wh portable power stations for field sensor deployments.\cite{power-stations} Table~\ref{tab:power} summarizes the evaluated models.

\begin{longtable}{>{\raggedright\arraybackslash}p{3.2cm} >{\raggedright\arraybackslash}p{1.8cm} >{\raggedright\arraybackslash}p{2.6cm} >{\raggedright\arraybackslash}p{2cm} >{\raggedright\arraybackslash}p{1.5cm} >{\raggedright\arraybackslash}p{2cm} >{\raggedright\arraybackslash}p{2.2cm}}
\caption{Portable power station comparison metrics.}\label{tab:power}\\
\toprule
Model & Capacity (Wh) & AC Output & USB-C Ports & Weight (lb) & Price (USD) & Notes \\
\midrule
GoLabs i200 & 256 & 200~W (250~W surge) & 60~W + 30~W & 8.37 & 229.99 & Reference unit; discontinued. \\
Jackery Explorer 240 v2 & 256 & 300~W (600~W surge) & 2\,×\,100~W + 15~W & 7.94 & 199.99 & Widely stocked (Amazon, Walmart). \\
Jackery Explorer 300 & 293 & 300~W (500~W surge) & 60~W & 7.10 & 178.00 & Lowest price among tested units. \\
Jackery Explorer 300 Plus & 288 & 300~W (600~W surge) & 2\,×\,100~W & 8.20 & 269.00 & Current availability through Walmart. \\
Anker SOLIX C300 & 288 & 300~W (600~W surge) & 2\,×\,140~W + 15~W & 9.10 & 229.99 & 50-minute fast charge reported. \\
EcoFlow RIVER 2 & 256 & 300~W (600~W surge) & 60~W & 7.70 & 169.99 & 60-minute recharge; dual AC receptacles. \\
BLUETTI EB3A & 268.8 & 600~W (1.2~kW surge) & 100~W & 10.10 & 189.99 & Highest AC rating; heavier chassis. \\
Goal Zero Yeti 300 & 297 & 350~W (600~W surge) & 100~W + 2\,×\,30~W & 13.70 & 349.95 & Premium pricing; limited stock outside Amazon. \\
\bottomrule
\end{longtable}

\section*{Conclusion}
Across sensing, data curation, hardware readiness, and algorithm development, the archived artifacts document a mature research stack ready for expanded fielding and additional automation. Continued work can build directly on the delivered calibration assets, fused datasets, HDR capture workflows, and benchmarking toolchains captured here.

\begin{thebibliography}{99}
\bibitem{gpu-summary-table} J.~C.~Vaught, ``GPU Parallelism Literature Review Summary Table,'' 2025.
\bibitem{gpu-weekly-brief} J.~C.~Vaught, ``Radar Update Weekly Brief -- 23 September 2025,'' 2025.
\bibitem{thermal-summary} J.~C.~Vaught, ``Thermal Calibration Summary,'' 2025.
\bibitem{maritime-datasets} J.~C.~Vaught, ``Overview of Maritime Datasets for Vision-Based Navigation and Detection,'' 2025.
\bibitem{fusion-summary} J.~C.~Vaught, ``Radar-Camera Fusion Summary,'' 2025.
\bibitem{monthly-report-may2025} University of South Carolina \\ Integer Technologies, ``Advanced Perception for Autonomous Platforms in the Littorals: Monthly Scientific and Technical Report 22,'' May 2025.
\bibitem{phase2-gain} J.~C.~Vaught, ``Phase 2 Gain Analysis Report,'' June 2025.
\bibitem{phase2-hdr} J.~C.~Vaught, ``Phase 2 HDR Tracking Report,'' July 2025.
\bibitem{lake-access} J.~C.~Vaught, ``SC Lake Radar Access Locations Workbook,'' 2025.
\bibitem{power-stations} J.~C.~Vaught, ``Portable Power Station Comparison Workbook,'' 2025.
\end{thebibliography}

\end{document}
